% dissemination_titlepage.tex

\vspace*{2.5cm}

\begin{center}
  {\LARGE \textbf{Dissemination Strategy Guide}\par}
  \vspace{1.0cm}
  {\large \textbf{\#TERRA: Paths to Sustainability in Adolescence (COMPETE2030-FEDER-00903800)}\par}
  \vspace{1.5cm}
\end{center}

\begin{center}
\textbf{Authors}

\vspace{0.5cm}

Tiago Ferreira\\
\href{https://orcid.org/0000-0002-2884-2547}{ORCID: 0000-0002-2884-2547}\\
Center for Psychology at the University of Porto (CPUP)\\
Faculty of Psychology and Education Sciences, University of Porto

\vspace{0.6cm}

Filipa Nunes\\
\href{https://orcid.org/0000-0002-1315-3327}{ORCID: 0000-0002-1315-3327}\\
Center for Psychology at the University of Porto (CPUP)\\
Faculty of Psychology and Education Sciences, University of Porto

\vspace{0.6cm}

Mar\'{i}lia Montenegro\\
Center for Psychology at the University of Porto (CPUP)\\
Faculty of Psychology and Education Sciences, University of Porto

\vspace{1.2cm}

{\large \today}
\end{center}

\vspace{1.5cm}

\noindent
\textbf{Summary.} This Dissemination Strategy Guide outlines the communication and dissemination approach for the \#TERRA project, covering both internal and external channels. It presents the project's visual identity, defines the goals of the communication strategy, and describes the dissemination tools and activities planned throughout the project lifecycle (September 2025 -- September 2028). Special attention is given to promoting project results and ensuring transparent, timely, and inclusive information flow among the research team, advisory panel, stakeholders, and the broader community.

\vspace{0.5cm}

\noindent\textit{The content of this document represents the authors' views, being their sole responsibility.}

\vfill
